% ============================================================================
% ABSTRACT + KEYWORDS - Según Anexo 8 de la Pauta UTEM
% ============================================================================
% Versión en inglés del resumen y las palabras claves.
% ============================================================================

\chapter*{ABSTRACT}
\thispagestyle{fancy}

This study addresses the logistical and operational inefficiency in vehicular access control at the Ñuñoa Campus of the Universidad Tecnológica Metropolitana (UTEM). Currently, staff entry depends on manual intervention by security personnel, leading to prolonged waiting times and diverting human resources from critical institutional security tasks. To resolve this issue, the overall objective of the project is to validate the technical and operational feasibility of an automated access control system by developing a functional prototype based on Artificial Intelligence (AI) for the recognition of Chilean license plates.

The development methodology is based on a hybrid approach that combines the Scrum agile framework for project phase management with the CRISP-DM standard to govern the deep learning model's lifecycle. Architecturally, a solution was designed under the Edge Computing paradigm (Technological Retrofit), shifting the inference of the YOLOv12s neural network directly to the edge node (security gate) to guarantee response times under 100 ms and ensure system operation during network outages (Offline-First).

In this initial stage of research and design, the results demonstrate the structural viability of the technical proposal and its high institutional impact. The economic feasibility study confirms that modernizing the access via edge infrastructure is a highly profitable investment for the university, achieving a Net Present Value (NPV) of \$5,898,595 CLP, an Internal Rate of Return (IRR) of 59.47\%, and a rapid Payback Period of 1.55 years. This theoretical and practical design establishes a robust foundation for the software development, empirical model training, and technological integration testing phases corresponding to the next stage of the project.

\noindent\textbf{KEYWORDS:}\\
Computer Vision, Artificial Intelligence, Edge Computing, CRISP-DM, Project Evaluation.

\newpage
